%\documentclass[draft]{beamer}
\documentclass[xcolor=pdftex,svgnames,table]{beamer}
%\usepackage[table]{xcolor}
\usepackage{latexsym,amssymb,amsmath,stmaryrd}
\usepackage[utf8]{inputenc}
\usepackage[T1]{fontenc}
\usepackage[francais]{babel}
\usepackage{pgf}
\usepackage{tikz}
\usetikzlibrary{arrows,patterns,plotmarks,shapes,snakes,er,3d,automata,%
backgrounds,topaths,trees,petri,mindmap}
\usepackage{pgfbaseimage}
\usepackage{ifpdf}
\ifpdf
\usepackage{graphicx}
\else
%\usepackage[dvips]{graphicx}
\usepackage{pstricks,pst-tree,pst-node}
\newpsobject{showgrid}{psgrid}{subgriddiv=1,griddots=10,gridlabels=6pt}
\fi
\usepackage{moreverb}
\usepackage[normalem]{ulem}
\usepackage{version}
\usepackage{url}
\usepackage{multicol}
\usepackage{listings}
\lstset{
  language=[ANSI]C, 
%  gobble=2, 
%  escapeinside="", 
  basicstyle=\ttfamily, 
%  directivestyle=\color{Fuchsia},
  identifierstyle = \color{DarkOrange},
  keywordstyle=\color{DarkGreen}, 
  commentstyle=\color{red}, 
%  numbers=left, 
%  numbersep=5pt, 
%  numberstyle=\scriptsize, 
%  moredelim=[is][\color{gray}\itshape]{/*}{*/}, 
%  moredelim=[is][\alert]{/+}{+/}, 
%  morecomment=[is]{/=}{=/}, 
%  morecomment=[in]{comment=}{,} 
%  fancyvrb=true 
} 

\newcommand{\binaire}[1]{\ensuremath{\underline{#1}}}
\newcommand{\C}[1]{\texttt{#1}}
\newcommand{\bbbn}{\ensuremath{\mathbb{N}}}
\input{communs/cc_beamer}
\usetheme{classic}
\newcommand{\nowrite}{\put(10,-4){\includegraphics[scale=.05]{creative_commons/nopencil}}}
\newcommand{\youwrite}{\put(10,-4){\includegraphics[scale=.05]{creative_commons/pencil}}}
\newcommand{\aemporter}{\put(10,-6){\includegraphics[scale=.05]{creative_commons/szymonraj_Shopping_bag}}}



\title{Bases de programmation -- Cours 6.\\ Jeu d'adresses.}
\author{Pierre Boudes}
\date{\today}

\begin{document}

%% Page de titre et licence CC.
\begin{frame}
        \titlepage
        \vfill
        \begin{center}
                \CcGroupByNcSa{0.83}{0.95ex}\\[2.5ex]
                {\tiny\CcNote{\CcLongnameByNcSa}}
                \vspace*{-2.5ex}
        \end{center}
\end{frame}

\section[Plan]{}
\frame[label=plan]{\tableofcontents}

\section{Adresses mémoires, pointeurs}

\begin{frame}[fragile]
  \frametitle{Adresses mémoires, pointeurs}


  \begin{itemize}
  \item  Si \C{x} est une variable, \C{\&x} est son adresse mémoire.
  \item  Si \C{p} contient une adresse mémoire, \C{*p} est la valeur contenue
    à cette adresse. Et on dit que \C{p} est un \alert{pointeur} car son rôle
    est de désigner (pointer en direction de) une autre case mémoire.
  \end{itemize}
\begin{lstlisting}[numbers=left,basicstyle=\ttfamily\small]
int main () {
  int a = 3;
  int b = 4;
  int *p; /* /!\ declaration d'un pointeur */
  p = &a; /* p contient l'adresse de a */
  *p = 1; /* ecrit a l'adresse contenue dans p */
  p = &b;
  *p = 2;
  printf("%d %d\n", a, b);
  return EXIT_SUCCESS;
}
\end{lstlisting}
\end{frame}

\begin{frame}[fragile]
  \frametitle{Indirection}

\begin{columns}
\begin{column}[b]{3.7cm}

\begin{lstlisting}[basicstyle=\ttfamily]
  int a = 5;
  int b = 4;
  int *p = &b;
  a = 3 * a + *p;
\end{lstlisting}

\C{
\begin{tikzpicture} [level distance=9mm,every node/.style={fill=blue!20,circle,inner sep=1pt,minimum size=18pt}]
\node {=}
child {node {a}}
child {node {+}
  child {node {*}
    child {node {3}}
    child {node {a}}
  }
  child {node {*p}}
};
\end{tikzpicture}
}
\end{column}
\begin{column}[b]{3.7cm}
\begin{lstlisting}[numbers=left,basicstyle=\ttfamily\small]
valeur 3 r0
lecture 10 r1
mult r0 r1
lecture 12 r2
lecture *r2 r3
add r1 r3
ecriture r3 10
stop

5  # a
4  # b
11 # p
\end{lstlisting}
\end{column}
\end{columns}
\end{frame}

\begin{frame}
  \frametitle{Erreurs de casting}
erreurs makes pointer form integer without a cast.
\end{frame}

\begin{frame}[fragile]
  \frametitle{Chaînes}

Pour les chaînes \C{char * s} plutôt que \C{char s[]}.

\begin{lstlisting}[numbers=left,basicstyle=\ttfamily\small]
#include <stdio.h>

typedef char mstring[];
typedef char *istring;
int main () {
  mstring toto = "Salut mon Toto";
  istring tata = "Bonjour Tata";
  char * personne;
  personne = tata;
  printf("%s\n", personne);
  // personne[5] = '\0'; <-- bus error 10
  personne = toto;

  printf("%s\n", personne);
  personne[5] = '\0';
  printf("%s\n", personne);
  printf("%s\n", toto);
  return 0;
}

\end{lstlisting}


\end{frame}

\section{Adresses et fonctions}
\subsection{Argument-résultat}
\begin{frame}[fragile]
  \frametitle{Argument résultat}

\begin{lstlisting}[numbers=left,basicstyle=\ttfamily\small]
int age = 0;
scanf("%d", &age);
\end{lstlisting}
\end{frame}


\begin{frame}[fragile]
  \frametitle{Focus sur le passage de valeurs}
  Les fonctions communiquent \textbf{des valeurs}, pas des noms de
  variables.

\begin{columns}
\scriptsize
  \begin{column}[t]{3.7cm}
 \begin{lstlisting}[numbers=left,basicstyle=\ttfamily\scriptsize]
void permute_valeurs(int a,int b);

int main() {
    int x = 1;
    int y = 2;
    permute_valeurs(x, y);
    printf("x = %d et y = %d\n",x,y);
    return EXIT_SUCCESS;
}

void permute_valeurs(int a, int b) {
    int aux;
    aux = a;
    a = b;
    b = aux;
}\end{lstlisting}
\vspace{.4cm}
 \end{column}
\end{columns}
\end{frame}



\begin{frame}[fragile]
  \frametitle{Focus sur le passage de valeurs}
  Les fonctions communiquent \textbf{des valeurs}, pas des noms de
  variables.
\begin{columns}
\scriptsize
\begin{column}[t]{3.7cm}
 \begin{lstlisting}[numbers=left,basicstyle=\ttfamily\scriptsize]
void permute_valeurs(int a,int b);

int main() {
    int x = 1;
    int y = 2;
    permute_valeurs(&x, &y);
    printf("x = %d et y = %d\n",x,y);
    return EXIT_SUCCESS;
}

void permute_valeurs(int *a, int *b) {
    int aux;
    aux = *a;
    *a = *b;
    *b = aux;
}\end{lstlisting}
\vspace{.4cm}
 \end{column}
\end{columns}
\end{frame}


\subsection{Passage de tableau en paramètre}
\begin{frame}
  \frametitle{Passage de tableaux en paramètre}
  Lorsqu'on passe un tableau en paramètre d'une fonction le paramètre
  formel qui le reçoit doit être un pointeur (du bon type).
\end{frame}



\subsection{Passage de structures par adresse (efficacité)}
\begin{frame}[fragile]
  \frametitle{Passage de structure par adresse}
  Pour rendre les appels plus rapides on passe souvent juste l'adresse
  d'une structure en paramètre d'une fonction plutôt que de passer
  toute la structure. Cela évite la copie de la structure.

  Nouvelle notation (pratique) \verb+->+ : si \verb+fiche+ est un pointeur
  sur une structure, plutôt que d'écrire \verb+(*fiche).nom+ on écrira \verb+fiche->nom+.
\end{frame}

\section{Allocation mémoire}
\begin{frame}
  \frametitle{Allocation mémoire}
  Pouquoi vous voudrez apprendre malloc et free.
\end{frame}

% \subsection{Pour aller plus loin en C}
% \begin{frame}
%   \frametitle{Pour aller plus loin\nowrite}

% \begin{itemize}
% \item La taille d'un tableau statique gagne à être fixée par une constante symbolique (\C{\#define N 3}).
% \pause
% \item L'identificateur du tableau (\emph{ie} \C{tab} dans la déclaration \C{int tab[3];}) est lui même une variable. Sa valeur est l'adresse de la première case du tableau \C{tab[0]}.
% \pause
%   \begin{itemize}
%   \item  Les variables dont la valeur est une adresse s'appellent des \alert{pointeurs};
% \item La notation \emph{esperluette}, \C{\&x}, donne accès à l'adresse d'une variable;
% \item La notation \emph{étoile}, \C{*tab}, ne s'applique qu'à une adresse, elle donne alors accès à la valeur contenue à cette adresse.
% \item Les expressions \C{tab[i]} et  \C{*(tab + i)} sont identiques en
%   C.
%   \end{itemize}\pause
% \item Pour les grands tableaux, il est nécessaire d'allouer dynamiquement l'espace mémoire.\pause
% \item Si on ne veut pas s'ennuyer avec les pointeurs (et l'allocation)
%   on peut emballer nos tableaux statiques dans des \structure{struct}.

% \end{itemize}
% \end{frame}


% \section{Conclusion}
% \begin{frame}[fragile]
%   \frametitle{Conclusion}
%   \begin{itemize}
% \item On évitera de passer des tableaux à des fonctions (pointeur, mémoire), on les
%   emballera avant dans des structures :
%   \begin{lstlisting}
% #define TAILLE_MAX 1024
% typedef struct tableau_s {
%     int taille;  /* maximum TAILLE_MAX */
%     int t[TAILLE_MAX];
% } tableau_t;
%   \end{lstlisting}
%   \end{itemize}
% \end{frame}

\end{document}
